\documentclass[11pt]{article}
\pdfinfoomitdate=1\pdftrailerid{}
\usepackage[letterpaper,margin=.7in,headheight=15pt,headsep=16pt,footskip=24pt]{geometry}
\usepackage[T1]{fontenc}
\usepackage{lmodern,amsmath,amssymb,array,booktabs,fancyhdr,hyperref}
\renewcommand{\familydefault}{\sfdefault}
\setlength{\parindent}{0pt}\setlength{\parskip}{6pt}
\hypersetup{hidelinks,pdftitle={Week 63 Cards away from home Adult return-visit guide}}
\pagestyle{fancy}\fancyhf{}
\fancyhead[L]{\small Week 63 / Cards away from home / Adult return-visit guide}
\fancyfoot[L]{\scriptsize Bellingham Math Circle / Week 63 / W63-RV-FAC-v2}
\fancyfoot[R]{\thepage}
\newcommand{\head}[1]{\par\textbf{\large #1}\par}
\newcommand{\sub}[1]{\par\textbf{#1}\par}
\newcommand{\row}[1]{\texttt{#1}}
\newcommand{\hint}[1]{\emph{If needed:} #1}
\begin{document}
\head{Mathematical overview}
\textbf{Fewest swaps equals cards minus loops.} For $n$ distinct cards in fixed, equally labeled homes, read rows in home order. Exchange \emph{any two distinct cards}, not just neighbors. Arrows run card-to-home, counting each home match as a one-card loop. With $c$ loops, exactly $n-c$ swaps bring the row home.

A swap splits one loop or joins two (worked examples on guide page 5). Each move changes the loop count by one; home has $n$ loops. For attainability, exchange a misplaced card with the card in its home, splitting off a singleton; repeat. Thus the hardest rows are single $n$-loops needing $n-1$ swaps. For $n\ge2$, a row can come home in exactly $m$ swaps if and only if $m\ge n-c$ and $m$ has the same parity as $n-c$: every move flips loop parity, and a repeated swap adds two moves without changing the finish.

\textbf{Leaving every home needs $\lceil n/2\rceil$ swaps, for $n\ge2$.} Each move touches at most two cards. For even $n$, exchange disjoint pairs. For odd $n\ge3$, pair all but three cards, then use two swaps on those three. This proves both necessity and sufficiency. There is no home-free row when $n=1$.

\textbf{A complete deck averages one home match per row.} Across all $n!$ rows, a chosen card is home in $(n-1)!$ rows. Summing by card gives $n(n-1)!=n!$ marks, for $n\ge1$. Pairs of home matches total $\binom n2(n-2)!=n!/2$, for $n\ge2$. These are totals over a complete deck, not guarantees about one row or an arbitrary sample. If $D_j$ counts home-free rows on $j$ labels, then exactly $k$ matches occur in
\[
\binom nk D_{n-k}\quad\hbox{rows},\qquad D_0=1,\quad D_1=0.
\]
Choose the matched labels; the remaining labels must all miss their own homes. Unlike the first visit's overlapping home-groups, these exact piles are disjoint. The zero-match and one-match piles differ by one: $D_n-nD_{n-1}=(-1)^n$. The first visit's recurrence proves the alternating difference; see guide page 6.

\textbf{Forbidden spots can overlap in different ways.} For an arbitrary board, let $r_k$ count choices of $k$ forbidden card--home spots using distinct cards and distinct homes, with $r_0=1$. Avoiding every forbidden spot leaves
\[
\sum_{k=0}^{n}(-1)^k r_k(n-k)!\quad\hbox{rows}.
\]
Compatible required spots leave $(n-k)!$ completions; clashing spots leave none. Inclusion--exclusion cancels every row that uses a forbidden spot and retains every legal row once. Equal numbers of forbidden spots need not give equal counts. The own-home recurrence does \emph{not} apply to arbitrary boards.

\textbf{Discoveries by band.} K--1: three or four cards can all leave home in two swaps; five need three because two swaps leave a card untouched. The A--C deck has as many home-match counters as outcomes, two per letter. Middle: the six hardest four-card rows need three swaps, and one swap homes two cards exactly when they occupy each other's homes. The deck's circles total its outcomes; equal numbers of crosses need not leave equal counts. ``No five'' and ``nine is most'' remain children's conjectures. Upper: cards minus loops gives the minimum and parity; totals and exact piles can be predicted without listing. The alternating $\pm1$ stays a conjecture without the recurrence. Strip/ring catalogs lead to repaired overlap counts.

New this visit: any-two-card swaps and disjoint exact piles; the first visit's home-groups overlapped. Only the upper table met the base page 7 arrows. Experiments suggest; complete catalogs settle finite hunts; arguments establish general facts.
\newpage
\head{Preparation and a first hour}
\textbf{Status.} Unpiloted and physically unrehearsed: printing, cutting, placement, counters and timings need a real trial. Print US Letter at 100\%, single-sided. Student pages: K--1, 5; middle, 6; upper, 10.

\textbf{Reuse W63-M-v3, \row{week-63-materials.pdf}.} Per pair: materials page 1 supplies two A--E decks (10 cards, $30\times40$ mm) and two strips of five $35\times45$ mm homes. Add cover slips, pencil, scrap and tray. Five kits serve two young pairs, two middle pairs and an upper trio rotating roles. A row uses one of each letter.

\textbf{Outcome work.} Per young pair: one cut set of six A--C outcome cards from materials page 2 and a cup of 12--20 counters no wider than 12 mm; do not pre-count six. Five A--C sets exist in the base kit. Give the young table one extra uncut page 2 for a shared pencil alternative. Middle and upper tables each need one cut 24-card A--D deck. Each middle pair gets its own unmarked, uncut page 3; the upper trio shares one. These are \textbf{four extra uncut sheets}: one page 2 and three page 3. Never pass an already-circled sheet to a new pair for Problem 5.

Whole outcomes: $56\times27$ mm; cells: A--C $14\times10$ mm, A--D $9\times10$ mm. Supply pencils/erasers. Sort cut decks on the table; count distinct row IDs once. Without a retained kit, the visit's minimum is page 1 $\times5$ cut, the six A--C outcomes from page 2 $\times2$ cut, page 3 $\times2$ cut, plus the four uncut sheets above. Leave page-2 labels/blank records uncut. Cutting this minimum is unmeasured; the base full-kit forecast is an untested 45--90 minutes. Retained-kit checking, extras and counters: an untested 10--15 minutes.

\textbf{Which surfaces are mats.} Keep strips whole; align cards at the lower edge, leaving 5 mm visible. Cover D/E for three cards or E for four. Working mats: K--1 pp. 4--5, upper pp. 7--8 ($35\times45$ mm homes). Middle pp. 5--6 are reference boards, not mats. Small boxes/grids are pencil records. No base student packet is needed.

\textbf{Launch together, about 5 minutes after movement.} Let everyone handle cards, then point to the unchanged homes. Make four scrap-paper cards U, V, W, X (under a minute) and draw fixed U--X homes. Lay the cards as V U W X. Ask a child to physically trade V and X, near or far, while another names them; write \row{VX}. This differs from neighbor-only swaps. Record card letters, not home positions.
\begin{center}\small
\begin{tabular}{ccc}
Input in homes U V W X & Choose cards V and X & Output in the same homes\\[3pt]
\row{V U W X} & trade the first and last cards & \row{X U W V}\\
 & record \row{VX} once &
\end{tabular}
\end{center}
Count swaps afterwards; keep one start/final row rather than simultaneous tallies. Introduce arrows (upper p. 3) and grid dots (p. 9) at first use, with their printed non-task visuals.

\textbf{First route and staffing.} One adult anchors each table. K--1: picture matching, counting to ten; parent reads and scribes, children choose. Use 1--2, selected rows of 4, then 5; save 3 and 6--8. Middle: four labels, totals to 24; adult reads aloud and may record sums. Choose 1--4 then 5 or 7; second visit 6--9. For 9 each trial checks 24 outcomes and takes minutes: partners check half the deck each against one frozen board. Upper: five letters, swap words and totals to 120; adult reads and may scribe children's choices. Use 1 and selected rows of 2, then 3 or 5. Second visit: 4 and 6--9; third: 10--13, grids last. In 10 split branches by the card in home A and pool the 16 rows. In 12--13 share the 15 compatible-pair and 10 triple checks on the four-card strip; the larger grids need more, not faster bookkeeping.

For 60 minutes: 5 movement, 5 launch, 20 exploration, 2 stretch, 18 continuation, 5 sharing, 5 packing. Adults record rather than choose cases or moves; finishing every page is not the goal.
\newpage
\head{K--1 solutions and held hints}
\textbf{Parent's kit and route.} Each pair: W63-M-v3 p. 1 kit (two decks, two strips), D/E (or E) covered; cut A--C outcomes from p. 2; cup of 12--20 counters, at most 12 mm; pencil/scrap. Read/scribe; children choose. Start with 1--2, selected rows of 4, then 5; save 3 and 6--8.

Rows are in home order. Words such as \row{AB, BC} name the cards exchanged, wherever they now sit; the adult writes while the child acts.

\sub{Problem 1, student page 1}
The fewest is \textbf{2 swaps}; either \row{BCA} or \row{CAB} can result. One swap leaves the third card untouched and still home. For example, from \row{ABC}, swapping \row{AB} then \row{AC} gives \row{BCA}. Completeness: A uses B or C, forcing the rest. \hint{``Which card did that swap never touch?''} Let partners invent routes.

\sub{Problem 2, student page 2}
The fewest is \textbf{2}; all possible shortest finishes are \row{BADC}, \row{CDAB}, \row{DCBA}. Two swaps must touch all four cards in disjoint pairs. A partners B, C or D; the other two pair. This proves completeness. \hint{``Can one card take part twice if every card has to leave home?''}

\sub{Problem 3, student page 2}
\textbf{Two cannot; three can.} Two swaps touch at most four of the five cards. From \row{ABCDE}, the word \row{AB, CD, AE} gives \row{BEDCA}, with nobody home. \hint{``Is there a card that has never moved?''} Point to the card that never moved: two swaps are not enough.

\sub{Problem 4, student page 3}
\begin{center}\renewcommand{\arraystretch}{1.15}
\begin{tabular}{lcl}
Printed row & Fewest & One shortest word to home\\\midrule
\row{CBA} & 1 & \row{AC}\\
\row{BCA} & 2 & \row{AB, BC}\\
\row{BADC} & 2 & \row{AB, CD}\\
\row{BCDA} & 3 & \row{AB, BC, CD}
\end{tabular}
\end{center}
\row{CBA} is one exchange; \row{BADC} has two swapped pairs. \row{BCA} has no pair in each other's homes, so needs two. In \row{BCDA}, any first swap leaves three mixed cards or two swapped pairs; either needs two more. \hint{``Could this swap put both cards home? Which homes are they in now?''} Later: invent a row with a chosen minimum.

\sub{Problem 5, student page 4}
\textbf{Six counters, six outcomes, two counters for each letter.} In printed outcome order \row{ABC, ACB, BAC, BCA, CAB, CBA}, the match counts are $3,1,1,0,0,1$. A is home in \row{ABC, ACB}; B in \row{ABC, CBA}; C in \row{ABC, BAC}. \hint{``Can we line up one counter beside each whole outcome? Which outcomes give A a counter?''} If counters obscure matches, point to each match then tally it, or use the extra uncut sheet for shared pencil circles.

\sub{Problems 6--8, student pages 4--5}
\textbf{6:} only \row{BCA}. A cannot use A or B, so it must use C; then C must use B, leaving B in A. \textbf{7:} exactly \row{BCDA, CDAB, DCAB}. Split by home A: B gives only \row{BCDA}; C gives only \row{CDAB}; D gives only \row{DCAB}. \textbf{8:} only \row{CDAB}. Home B can take only D and home C only A. Then home A refuses A and B, so takes C; home D takes B: \row{CDAB}. These forced checks also prove completeness.

For board 6, point to home B and name both crossed-out pictures. Ask which cards may go there, then let children place them.

\hint{``Which home has only one allowed card?''} Children choose/check. Spare boxes do not promise counts. Point to unclear crosses with their home; handling is untested.
\newpage
\head{Grades 2--3 solutions and held hints}
\sub{Problem 1, student page 1}
Four cards need \textbf{2 swaps}, ending in \row{BADC, CDAB} or \row{DCBA}; disjoint pairs prove completeness (guide page 3). Five cards need \textbf{3}, for example \row{AB, CD, AE} from home gives \row{BEDCA}. Two moves cannot touch all five. \hint{``Which cards have not moved?''}

\sub{Problem 2, student page 2}
\begin{center}\small\renewcommand{\arraystretch}{1.1}
\begin{tabular}{lcl@{\qquad}lcl}
Row & Fewest & Word to home & Row & Fewest & Word to home\\\midrule
\row{CBAD}&1&\row{AC}&\row{BCAD}&2&\row{AB, BC}\\
\row{DCBA}&2&\row{AD, BC}&\row{ACDB}&2&\row{BC, CD}\\
\row{BCDA}&3&\row{AB, BC, CD}&\row{CDBA}&3&\row{AC, BD, CD}
\end{tabular}
\end{center}
One swap homes two cards only if they occupy each other's homes. \hint{``After putting one card home, what remains?''}

\sub{Problem 3, student page 2}
The six hardest rows: \row{BCDA, BDAC, CADB, CDBA, DABC, DCAB}, each needing \textbf{3 swaps}. Four never are needed: put A, then B, then C home, skipping matches; D is forced. A row with a card home needs at most two swaps on the other three. So hardest rows lie among the first visit's nine home-free rows. Three, \row{BADC, CDAB, DCBA}, are two swapped pairs needing two moves; the other six are those listed. \hint{``Which of the nine can you bring home in two?''} Wasted swaps do not count as needed.

\sub{Problem 4, student page 3}
\textbf{Yes:} \row{CDAB} (exchange \row{AC} or \row{BD}) and \row{DBCA} (\row{AD}). \textbf{No:} \row{BCDA, BCAD, DACB}. Two newly home cards must previously occupy each other's homes; conversely that exchange works. \hint{``Where is the card whose home this one occupies?''}

\sub{Problems 5--6, student pages 3--4}
\textbf{5:} 24 circles, six per letter. Fix A home: B,C,D have six orders; likewise each letter, giving $4\times6=24$. \textbf{6:} A--C gives six circles ($3\times2$); A--E gives 120 ($5\times24$). With A home, the other four arrange as many ways as the A--D deck has cards: 24. \hint{``How many rows keep this letter home, allowing others home too?''} The five-card prediction is optional abstraction.

\sub{Problem 7, student page 4}
For 0,1,2,3,4 home letters, the piles have \textbf{9,8,6,0,1} outcomes. Exactly one: choose its letter in four ways, then the remaining three must be home-free in two ways, giving $4\times2=8$. Exactly two: choose the two home letters in six ways and swap the other pair. Exactly three is impossible, since the last is forced home. \hint{``If those are the only home letters, what must the others avoid?''} Sort whole outcomes once, not once per match.

\sub{Problem 8, student page 5}
Top to bottom the counts are \textbf{9,7,4,0}. Filter the complete 24-card deck against each frozen board. On the third board A and B must use C and D, leaving C and D to use A and B: \row{CDAB, CDBA, DCAB, DCBA}. On the fourth, A has nowhere to go. \hint{``Does each cross remove a new group of rows, or can the groups overlap?''}

\sub{Problem 9, student page 6}
The maximum is \textbf{9}: cross each card out of its own home (any one-cross-per-card-and-home pattern works). For \textbf{8}, forbid A-in-A, B-in-B, D-in-C, D-in-D. For \textbf{6}, forbid A-in-A, B-in-B, D-in-A, D-in-B. Exactly \textbf{5 is impossible}; guide page 8 gives the adult proof. \hint{``Can your crosses share a card or home?''} Designs are checkable successes; a child's maximum remains a conjecture until explained.
\newpage
\head{Grades 4--5 swaps and loops}
\sub{Problem 1, student page 1}
The minima for three, four and five cards are \textbf{2,2,3}. The K--1 keys give constructions and complete shortest three/four-card finish lists. For five cards \row{AB, CD, AE} from home gives \row{BEDCA}; two swaps leave a card untouched. \hint{``How many different cards can two swaps touch?''}

\sub{Problems 2--3, student pages 2--3}
Each parenthesized group below is a directed loop, with the last letter pointing back to the first. A one-letter group is a self-loop. This is card-to-home notation, not a new student recording demand.
\begin{center}\small\renewcommand{\arraystretch}{1.2}
\begin{tabular}{llll}
Printed row & Loops & Fewest & One shortest word to home\\\midrule
\row{EBCDA}&$(AE)(B)(C)(D)$&1&\row{AE}\\
\row{BCADE}&$(ACB)(D)(E)$&2&\row{AB, BC}\\
\row{BADEC}&$(AB)(CED)$&3&\row{AB, CD, DE}\\
\row{BCAED}&$(ACB)(DE)$&3&\row{AB, BC, DE}\\
\row{BCDEA}&$(AEDCB)$&4&\row{AB, BC, CD, DE}\\
\row{DEABC}&$(ACEBD)$&4&\row{AD, BE, CD, DE}
\end{tabular}
\end{center}
\textbf{P2 hint:} ``Can this swap put a card home?'' \textbf{P3 hint:} ``How many separate loops are there before and after that swap?'' Use the printed \row{WUVXY} visual first: W sits in U, so $W\to U$, then $U\to V\to W$, with X and Y looping to themselves. In the task \row{BCADE}, A sits in C: $A\to C$, not $A\to B$.

On that non-task row, swapping W with V gives \row{VUWXY}: $U\to V\to W\to U$ splits into $U\leftrightarrow V$ and W home. Instead swapping X with Y gives \row{WUVYX}: the two one-card loops join into $X\leftrightarrow Y$. Cards trade their two arrow ends. In one loop this cuts it into two closed loops; in different loops it splices two into one. The child's rule: a loop of $k$ cards needs $k-1$ swaps; add these up, counting one-card loops too. Home has five separate loops, and a move gains at most one. The displayed words achieve $5-c$ moves. Inverting every arrow happens to preserve the loop count, but is still the wrong printed convention.

\sub{Problem 4, student page 4}
The hardest rows need \textbf{4 swaps}: exactly the single loops through all five labels. Start the loop at A; its other four labels can be ordered in $4!=24$ ways, with no rotational duplication. All 24 rows, grouped in alphabetical order for adult checking, are:
\begin{center}\ttfamily\small
\begin{tabular}{llllll}
BCDEA&BCEAD&BDAEC&BDECA&BEACD&BEDAC\\
CADEB&CAEBD&CDBEA&CDEAB&CEBAD&CEDBA\\
DABEC&DAECB&DCAEB&DCEBA&DEABC&DEBCA\\
EABCD&EADBC&ECABD&ECDAB&EDACB&EDBAC
\end{tabular}
\end{center}
Of the 44 home-free rows, the other \textbf{20 need 3 swaps}. Their loops have lengths two and three: choose the two labels in ten ways, and orient the three-loop in two ways, giving $10\times2=20$. No home-free five-card row has more than two loops. \hint{``If there are no one-card loops, which loop sizes can add to five?''} Children can count a family without copying this adult catalog.

\sub{Problem 5, student page 4}
For the requested exact lengths 2,3,4,5, \row{BACDE} gives \textbf{no, yes, no, yes}; \row{BCADE} gives \textbf{yes, no, yes, no}. Their shortest lengths are 1 and 2. Insert a repeated pair of swaps anywhere in a shortest word to add two. Loop parity rules out the other lengths, not merely a failed search. In general, for at least two cards, the possible lengths are the minimum plus $0,2,4,\ldots$. \hint{``What does doing the same exchange twice change? Can one exchange preserve loop parity?''} Detours are legal; a self-swap is not a move.
\newpage
\head{Grades 4--5 totals without listing}
\sub{Problem 6, student page 5}
The A--D sheet has \textbf{24 circles, six for each letter}. The predicted A--E deck has \textbf{120 circles}, 24 for each letter. Every outcome is counted once as an object, but can supply several marks. Fixing one card at home leaves respectively $3!=6$ or $4!=24$ arrangements. Sum over the four or five named cards. \hint{``Could we count A's circles first, without inspecting B's?''}

An average of one does not mean one on every row: the home row has four or five matches, while home-free rows have none.

\sub{Problem 7, student page 5}
There are \textbf{12 pairs} on A--D and \textbf{60 pairs} on A--E. A pair is an unordered choice of two circles on \emph{one} outcome. Four circles yield six pairs, not two. There are six choices of named letter-pair on four labels, and fixing that pair home leaves two orders of the remaining cards: $6\times2=12$. On five labels there are ten named pairs and six completions for each: $10\times6=60$.

For a concrete check, on \row{ABCD} point to the pairs AB, AC, AD, BC, BD, CD without removing the circles. The six exactly-two-match outcomes contribute one each; the all-home outcome contributes six, totaling 12. \hint{``How many outcomes have both A and B home, even if other letters are home too?''} Do not turn the question into pairing circles on different outcome cards. An optional later question is to count triples of matches; it is not needed for the first route.

\sub{Problem 8, student page 6}
\begin{center}\renewcommand{\arraystretch}{1.3}
\begin{tabular}{lrrrrrr}
Exactly this many matches & 0&1&2&3&4&5\\\midrule
A--D outcomes &9&8&6&0&1&--\\
A--E outcomes &44&45&20&10&0&1
\end{tabular}
\end{center}
For five cards: exactly one match means five choices of its label and nine home-free arrangements of the rest, $5\times9=45$. Exactly two means ten choices of matched pair and two home-free arrangements of the other three, $10\times2=20$. Exactly three means ten choices of matched triple and one exchange of the other two, $10\times1=10$. Exactly four is impossible; all five gives the one home row. The remaining pile has $120-45-20-10-0-1=44$ rows.

These subtractions need no overlap repair: each whole outcome has exactly one match count. Retain original labels and homes when considering the unmatched cards; do not compress them into a freshly relabeled task. \hint{``Which letters are exactly the home letters? What must happen to all the others?''} The first-visit facts $D_2=1,D_3=2,D_4=9$ are enough; there is no need to repeat its nine-row hunt.

\sub{Problem 9, student page 6}
The zero-match and one-match piles are \textbf{2 versus 3}, \textbf{9 versus 8}, \textbf{44 versus 45} for three, four and five cards. The difference has size one, switching which pile is larger. Trying these cases is evidence, not a proof that the pattern continues.

Here is a short adult explanation using the first visit's established recurrence, not reteaching its two-family construction. Put $\Delta_n=D_n-nD_{n-1}$. For $n\ge3$,
\[
\begin{aligned}
\Delta_n
&=(n-1)(D_{n-1}+D_{n-2})-nD_{n-1}\\
&=-\bigl(D_{n-1}-(n-1)D_{n-2}\bigr)=-\Delta_{n-1}.
\end{aligned}
\]
Since $\Delta_2=D_2-2D_1=1$, the difference is $(-1)^n$. \hint{``Can the exactly-one pile be counted by first choosing its one home letter?''} Offer the recurrence only once that count is clear. If the algebra would be adult-operated, keep the observed alternation as a conjecture for later explanation.
\newpage
\head{Grades 4--5 new forbidden homes}
\sub{Problem 10, student page 7}
The strip board forbids each card's own home and the home immediately to its right; the last card has no extra home at the right edge. Cover unused homes and set aside their cards, without relabeling. With A--C, only \row{BCA} obeys. With A--D, exactly \row{BCDA, CDAB, DCAB} obey. With all five cards, all \textbf{16} rows are:
\begin{center}\small
\begin{tabular}{llr}
Card in home A & Complete branch & Count\\\midrule
B & \row{BCDEA, BDAEC, BDEAC, BEDAC} & 4\\
C & \row{CDAEB, CDEAB, CDEBA, CEDAB, CEDBA} & 5\\
D & \row{DCAEB, DCEAB, DCEBA, DEABC} & 4\\
E & \row{ECDAB, ECDBA, EDABC} & 3
\end{tabular}
\end{center}
These first-card branches are disjoint and exhaustive. For each branch, place the remaining cards in home order, retaining every allowed choice and rejecting repeated cards; this tests all completions. Children can share branches and pool their records rather than each repeat the whole search. \hint{``Where can A still go? Which first-home cases have we not tried?''}

\sub{Problem 11, student page 8}
These boards add one prohibition: the last card cannot use home A. For four cards, remove \row{DCAB} from the preceding list, leaving exactly \row{BCDA, CDAB}. For five, remove \row{ECDAB, ECDBA, EDABC}, leaving these \textbf{13} rows:
\begin{center}\ttfamily\small
\begin{tabular}{lllll}
BCDEA&BDAEC&BDEAC&BEDAC&CDAEB\\
CDEAB&CDEBA&CEDAB&CEDBA&DCAEB\\
DCEAB&DCEBA&DEABC&&
\end{tabular}
\end{center}
Removal from a complete catalog proves completeness; it is not a new 120-row search. \hint{``Which rows of your old list fail the one new cross?''} The wraparound prohibition changes the allowed placements; it does not identify rotated or reflected rows.

\sub{Problem 12, student page 9}
The four-card grid answers, left to right, are \textbf{3 and 2}. Before counting, use the printed non-task \row{VXUW} bridge: V sits in U, so put its dot at card V, home U; the full output adds dots at X-in-V, U-in-W and W-in-X. A dot must not land on a shaded square. Let a child convert one already-built row to dots and back before asking for an overlap count. There is one dot per card-line and per home-column; axes must not be exchanged silently.

\sub{Problem 13, student page 10}
The five-card grid answers, left to right, are \textbf{16 and 13}. Use this only after the four-card count has meaning. The grids encode the same prohibitions as Problems 10--11, not additional constraints. The coefficient table below counts compatible \emph{forbidden} spots; it does not count legal full rows. Guide page 8 explains how to obtain it.
\begin{center}\small\renewcommand{\arraystretch}{1.15}
\begin{tabular}{lrrrrrrl}
Board & $r_0$&$r_1$&$r_2$&$r_3$&$r_4$&$r_5$ & Alternating count\\\midrule
4-card strip &1&7&15&10&1&--&$24-42+30-10+1=3$\\
4-card ring &1&8&20&16&2&--&$24-48+40-16+2=2$\\
5-card strip &1&9&28&35&15&1&$120-216+168-70+15-1=16$\\
5-card ring &1&10&35&50&25&2&$120-240+210-100+25-2=13$
\end{tabular}
\end{center}
\hint{P12: ``Can these two shaded spots happen in the same row?'' P13: ``How could we organize compatible choices without listing full rows?''} Use the spare printed grids for chosen spots, not a new simultaneous tally. A child may finish with the physical catalog and save inclusion--exclusion for another visit.
\newpage
\head{Explaining the counts and source trail}
\textbf{A workable overlap bridge for Problems 12--13.} Mark two forbidden spots sharing a card-line on a spare grid: no row can use both. Try two spots on distinct lines and columns: fixing them leaves two cards to place on a four-card board, hence two completions. Additional forbidden uses are allowed in this \emph{intersection} count. This is the new feature beyond first-visit own-home groups, where chosen matches never clash. Count compatible singletons, pairs and larger choices; use the table only after children have tried some.

For a systematic adult count, order the strip's forbidden spots as
\[
(A,A),(A,B),(B,B),(B,C),\ldots,(E,E).
\]
Only neighbors in this chain clash. Let $F(L,k)$ count ways to choose $k$ nonneighboring spots from a chain of length $L$. Its last spot is either absent or present, giving
\[
F(L,k)=F(L-1,k)+F(L-2,k-1),\qquad F(L,0)=1.
\]
Impossible choices count zero; $F(0,k>0)=0$, $F(1,1)=1$. This builds every strip coefficient with $L=2n-1$. For the ring of $L=2n$ spots, distinguish one spot: not chosen leaves a chain of $L-1$; chosen removes its two neighbors and leaves a chain of $L-3$. Thus $r_k=F(L-1,k)+F(L-3,k-1)$ for $k\ge1$, with $r_0=1$. These disjoint cases explain all coefficients without listing complete permutations.

The alternating signs can be justified on one frozen outcome. If it uses a forbidden spot, pair every subset of its forbidden uses that omits that spot with the subset that includes it. The two contributions have opposite signs and cancel. If it uses no forbidden spot, only the empty subset remains, counted once. This proves the formula for any finite forbidden board.

\textbf{Adult proof for middle Problem 9.} Four forbidden spots leave $24-4\cdot6+2r_2-r_3+r_4=2r_2-r_3+r_4$ rows. If all four use different cards and homes, $(r_2,r_3,r_4)=(6,4,1)$ gives 9. Otherwise $r_4=0$. If no compatible triple exists, the count $2r_2$ is even and at most 8: five compatible pairs among four spots would force a compatible triple. If a compatible triple exists, the fourth spot clashes with one or two of it (at most one sharing its card and one sharing its home). One clash gives $(r_2,r_3)=(5,2)$ and 8 rows; two give $(4,1)$ and 7. Thus a count is 9, 8, 7 or even and at most 8. Five is impossible and nothing exceeds 9. This is adult background, not a child's proof obligation.

\sub{Sources and consultation boundaries}
\small
\textbf{Writer consultation record (retained).} Final return-visit PDFs: K--1 pp. 1--5, P1--8; middle pp. 1--6, P1--9; upper pp. 1--10, P1--13, plus their sources, outline and both earlier reviews. Base Week 63 guide W63-G-v4, especially pp. 1--4 and 7--10, and materials W63-M-v3 pp. 1--3 establish the kit and prior mathematics. Style models: Week 8 return-visit guide pp. 1--5 and Week 98 guide, including its editable source.

Alan Frieze, CMU Discrete Mathematics, \href{https://www.math.cmu.edu/~af1p/Teaching/DM/D14.pdf}{D14 pp. 2--4} (derangement recurrence) and \href{https://www.math.cmu.edu/~af1p/Teaching/DM/D16.pdf}{D16 pp. 1--2, 4--7} (inclusion--exclusion and cancellation), are cited by the base source notes. Their originals were not reopened during this revision; no direct consultation is claimed here. The proofs and kit instances here were independently checked. The base source notes report \emph{Math Circle by the Bay}, preface printed pp. viii--x, and Rozhkovskaya, Lessons 3 item 1 and 8 item 2, for concrete attempts before records and reduced copying. Those books are unavailable in this checkout; that pedagogical evidence was consulted through the notes, not the originals. Our timing and route choices are untested adaptations.

\textbf{Not directly consulted.} The outline names Cayley (1849), Diaconis (1988), Stanley (2nd ed., \S1.3, \S2.2), Feller (vol. 1, ch. IV), Riordan (1958), Kaplansky (1943), Kaplansky--Riordan (1946), Touchard (1934), and algebra texts for sign. These are source leads, not claims of direct consultation. Record actual fit, preparation time, rule retention and which investigations children wish to continue. No six-card kit or classroom success is claimed.
\end{document}
